\section{Introduction}


Recent advances in large language models have enabled rapid progress on repository-level software engineering tasks, such as resolving real-world GitHub issues in large codebases~\citep{jimenez2024swebenchlanguagemodelsresolve}.

A critical prerequisite is \textbf{code localization}: given an issue description and a repository, an agent must identify the relevant files and finer-grained code entities (e.g., classes and functions) to modify~\citep{xia2024agentlessdemystifyingllmbasedsoftware}.

This is challenging in large repositories with complex interdependencies, where early agent steps are often dominated by navigation and search~\citep{cognition2025swegrep}. Inaccurate search wastes limited context budget and exacerbates performance degradation at long input lengths~\citep{hong2025context}.

To address this inefficiency, the community has begun exploring specialized \textbf{localization} subagents that identify relevant code before handing off to a downstream code editing agent, rather than relying on a monolithic agent to perform both search and editing~\citep{zhang2026toolenoughreinforcementlearning}.

Despite these advances, most high-performing code localization systems remain closed-source. Some open-source alternatives (e.g., LocAgent~\citep{chen-etal-2025-locagent}) use proprietary API trajectories (e.g., Claude-3.5) for distillation, which can limit reproducibility and deployment in privacy-sensitive settings. RepoNavigator~\citep{zhang2026toolenoughreinforcementlearning} demonstrates RL training without distillation for repository-level agents; in contrast, we focus on training a dedicated localization subagent with a lightweight tool scaffold and multi-level F1 rewards.

In this work, we present \textbf{an open-source recipe for training dedicated code localization agents using reinforcement learning}. Our contributions are three-fold: \yccomment{To be refined later.}
\sonicomment{We should add some summary of the results here saying that our training recipe is outperforming closed-source models + complex scaffolds, and even 8x larger models on Pro/Verified/Lite}
\begin{enumerate}
    \item \textbf{Multi-level localization with RL training}: We formulate code localization as multi-level target prediction at three granularities (file, class, and function), and optimize multi-level F1 rewards. We train \methodname with GSPO directly from a pretrained open-source model, without closed-source distillation, and use a simple turn-completion reward to stabilize multi-turn training.
    
    \item \textbf{Simple, reproducible agent scaffold}: We build \methodname on top of OpenHands with a minimal tool set (\texttt{Terminal} and \texttt{Localization\-Finish}). This avoids heavyweight repository indexing or graph construction, and does not require installing repository dependencies or containerized execution, while supporting multi-turn localization.
    
    \item \textbf{Competitive performance with full release}: \methodname achieves competitive localization performance on SWE-Bench Lite, SWE-Bench Verified, and SWE-Bench Pro~\citep{jimenez2024swebenchlanguagemodelsresolve,chowdhury2024swebenchverified,deng2025swebenchproaiagents}. We release training code, reward functions, model checkpoints, and evaluation datasets. \sonicomment{Do we actually need to submit the anonymized code here or just promise?}
\end{enumerate}

Our work demonstrates that with careful reward design and training techniques, open-source models can achieve competitive code localization performance without relying on expensive proprietary models or closed-source distillation.